\section{Uniform moment prices after changing Gaussian coordinates}
\label{sec:gaussian-moment-prices}
We now estimate the polynomial in the Gaussian sample introduced
by differentiation of the projected density. The covariance may
couple every coordinate to every other coordinate. Only its
uniform diagonal bound is needed for this step.

\subsection{A polynomial moment bound without independence}
Let $Z$ be a centered Gaussian vector with covariance $C\ge0$ and
$C_{ii}\le b$. Each coordinate has exponential moment
\[
 \mathbb E e^{tZ_i}=e^{C_{ii}t^2/2}\le e^{bt^2/2}.
\]
This follows directly by completing the square in its one-dimensional
Gaussian law, including the zero-variance case. For an integer
$r\ge0$,
\[
 (1+|z|)^r\le e^{r|z|}\le e^{rz}+e^{-rz}.
\]
Hence
\begin{equation}\label{eq:moment-coordinate}
 \mathbb E(1+|Z_i|)^r\le2e^{br^2/2}.
\end{equation}

Let $I$ be a set of physical bonds, let $I_{\rm col}$ contain all
their colour coordinates, and put $n=|I|$, $N=|I_{\rm col}|$.
Define
\[
 F_I(z)=1+n+N+\sum_{i\in I_{\rm col}}|z_i|,
 \qquad Y_I=1+n+2N.
\]
Then
\begin{equation}\label{eq:moment-face}
 \mathbb E F_I(Z)^r\le
 Y_I^r\{1+2Ne^{br^2/2}\}.
\end{equation}
For nonempty $I_{\rm col}$ choose a coordinate at which $|z_i|$
is maximal. The definition gives $F_I(z)\le Y_I(1+|z_i|)$, and
therefore
\[
 F_I(z)^r\le Y_I^r
       \left(1+\sum_{i\in I_{\rm col}}(1+|z_i|)^r\right).
\]
For an empty coordinate set the same inequality is immediate.
Take expectations and use \eqref{eq:moment-coordinate}. In
particular all these moments are integrable. This argument uses
one-coordinate marginals and a finite sum, without factoring a
joint moment.

For the local covariance, the weighted matrix bounds of
Section~\ref{sec:local-covariance} supply a fixed $b$, independent
of the period. Suppose an activity has nonnegative support price
$s$, equal to the sum of its covering numbers and nonnegative
length charges. A cover by radius-$B$ balls gives
\[
 |\operatorname{supp}|\le s\,b_B,\qquad
 |I|\le s\,b_B\rho_L,\qquad N=q|I|,
\]
where $q$ is the colour dimension and $\rho_L$ the fixed bond
density. The first inequality follows by summing the covering
balls over all factors; overlaps can only decrease the union.
The second uses that every bond of $I$ has its anchor in that union.
Since $b_B$ has upper exponent zero, \eqref{eq:moment-face} implies,
for every fixed $r$,
\begin{equation}\label{eq:moment-support}
 \mathbb E F_I(Z)^r\le A_r(x)(1+s)^{m_r},
\end{equation}
for an integer $m_r$ and a nonnegative $A_r$ of upper exponent
zero. Indeed insert the displayed bounds for $n,N$ into
\eqref{eq:moment-face} and expand the resulting fixed polynomial.

\subsection{The derivative of the density and Jacobian}
At the base parameter the compensated transformation is the identity.
The coordinate derivative bound following from
\eqref{eq:gaussian-compensated-displacement} is
\[
 \|D^r(T_yz)_i|_{y=y_0}\|
 \le\sqrt\beta\{u(x)+2m(x)\}+F_I(z),\qquad0\le r\le j,
\]
where $u,m$ bound respectively the compensation increment and mean
derivatives and have upper exponent $-1/2$. For the value, use
$|z_i|\le F_I(z)$; for positive orders differentiate the
compensated displacement. As $\beta\le2h$, the first term has
upper exponent zero.

The principal covariance and its inverse inherit the weighted
derivative bounds proved earlier. The finite-order relative
determinant argument of Section~\ref{sec:gaussian-differentiation}
therefore gives a derivative scale polynomial in $N$,
$F_I(z)$ and functions of upper exponent zero. The Jacobian is
a product of $n$ one-bond factors. Each has value one at the base
parameter and derivative scale of upper exponent zero, by
Section~\ref{sec:compensated-variation}. Its product has amplitude
one and the sum of these scales. Since $1,n,N,|z_i|\le F_I(z)$,
every term of the fixed polynomial is bounded by one common
power of $F_I(z)$. We obtain an integer $k$ and a nonnegative
$D(x)$ of upper exponent zero such that the complete density
ratio times Jacobian, denoted $R_y(z)$, satisfies
\begin{equation}\label{eq:moment-ratio-jets}
 R_{\cdot}(z)\in\mathcal B_j(1,D(x)F_I(z)^k;y_0).
\end{equation}
The exponent $k$ depends on the derivative order and fixed
construction data, not on the support size or period.

\subsection{Retaining the warning gain under integration}
Consider a finite product of local activity factors indexed by
$\mathcal P$. Write $t_i\ge0$ for their pointwise prices and
$s$ for their total support price. Suppose their proved local
derivative estimates, under the coherent variation, give
\begin{equation}\label{eq:moment-activity-premise}
 \Phi_{\cdot}(z)\in
 \mathcal B_j\left(\mathbf1_{\mathcal E}(z)
                   \prod_{i\in\mathcal P}t_i,\,2r(x)s;y_0\right),
 \qquad r\ge1,
\end{equation}
where $r$ has upper exponent zero and $\mathcal E$ is the joint
warning event for the $m$ warning factors. The displayed estimate
implies that all derivatives through order $j$ vanish outside
$\mathcal E$. In this application the stronger local-vanishing
property holds: an unsatisfied warning factor is zero on a
neighbourhood of $y_0$, as in the closed-support argument for
packet masks. For the
other factors, their derivative scales add: each nonzero supported
factor has covering number at least one, so its unit scale and
length scale together cost at most twice its assigned support
price. This proves the factor $2s$ in
\eqref{eq:moment-activity-premise} from the individual bounds.
The individual estimates for all species remain necessary inputs
to this product step.

Multiply \eqref{eq:moment-activity-premise} by
\eqref{eq:moment-ratio-jets}. Since $F_I\ge1$ and $1+s\ge1$,
\[
 1+DF_I^k+2rs\le(1+D+2r)(1+s)F_I^k.
\]
Converting to derivative scale one therefore bounds all derivatives
through order $j$ of the transformed integrand by
\begin{equation}\label{eq:moment-integrand-price}
 \mathbf1_{\mathcal E}(z)
 \left(\prod_{i\in\mathcal P}t_i\right)
 (1+D+2r)^j(1+s)^jF_I(z)^{kj}.
\end{equation}

Suppose the simultaneous-warning estimate of
Section~\ref{sec:gaussian-warnings} gives
$\mathbb P(\mathcal E)\le\rho^{2m}$ with $\rho\ge0$.
By Cauchy--Schwarz and \eqref{eq:moment-support} at order $2kj$,
\[
 \mathbb E\bigl[\mathbf1_{\mathcal E}F_I^{kj}\bigr]
 \le\mathbb P(\mathcal E)^{1/2}
             (\mathbb EF_I^{2kj})^{1/2}
 \le E(x)(1+s)^a\rho^m,
\]
for a nonnegative upper-exponent-zero function $E$ and an integer
$a$. For example replace the square root of the moment bound by
$(1+A_{2kj})(1+s)^{m_{2kj}}$. This also proves integrability of
\eqref{eq:moment-integrand-price}. Integrating any frozen Haar
variables adds no factor, since their total mass is one and the
estimate is uniform in them.

Consequently the integrated derivatives obey
\begin{equation}\label{eq:moment-final-price}
 \left\|D^r\int R_y\Phi_y\,d\mu\bigg|_{y=y_0}\right\|
 \le E(x)(1+D+2r(x))^j(1+s)^{a+j}
               \rho^m\prod_{i\in\mathcal P}t_i,
 \qquad 0\le r\le j.
\end{equation}
The derivative interchange is justified by the separate common
neighbourhood bound of
Section~\ref{sec:gaussian-differentiation}. Absorb one factor
$\rho$ into the price of each warning factor. The right-hand side
then has exactly the product-price form used by the connected
support summation, with one fixed polynomial in $1+s$. The
support-price estimate of Section~\ref{sec:mixed-activity-prices}
allows such a fixed polynomial through its exponential reserve.
Thus Gaussian differentiation preserves the available warning
gain while adding only upper-exponent-zero and polynomial support
costs. No independence of warning events or Gaussian coordinates
has been used.
